\section{Character projectors and the operator \texorpdfstring{$\Delta$}{Delta}}
\label{sec:character_projectors}

Throughout, $U$ is a finite-dimensional complex representation of a finite
group $G$.

\begin{lemma}[Group averaging]%
    \label{thm:peps_group_averaging}
    \lean{Representation.averageMap_linHom_apply,
        Representation.mem_invariants_linHom_iff,
        Representation.isProj_averageMap_linHom,
        Representation.averageMap_linHom_apply_of_adjoint,
        Representation.trace_adjoint_comp_averageMap_linHom}
    \leanok
    The map $\sigma(X)=\frac{1}{|G|}\sum_g U_gXU_g^{-1}$ is a projection onto
    the operators commuting with every $U_g$. If the representation is unitary,
    then $\sigma(X)=\frac{1}{|G|}\sum_g U_gXU_g^\dagger$ and
    $\tr[Y^\dagger\sigma(X)]=\tr[\sigma(Y)^\dagger X]$, so $\sigma$ is the
    orthogonal projector onto these operators for the Hilbert--Schmidt inner
    product \cite[Lemma~4.3]{Schuch2010PEPS}.
\end{lemma}

\begin{proof}\leanok
    The map $\sigma$ is the average of the conjugation action
    $X\mapsto U_gXU_g^{-1}$, so it is a projection onto the operators fixed
    by every conjugation, and these are the operators commuting with every
    $U_g$. For a unitary representation,
    $\tr[Y^\dagger U_gXU_g^\dagger]=\tr[(U_g^\dagger YU_g)^\dagger X]$, and
    summing over $g$, with $g$ replaced by $g^{-1}$, gives the symmetry of
    $\sigma$.
\end{proof}

\begin{definition}[Isotypic projectors and the operator $\Delta$]%
    \label{def:peps_delta_operator}
    \lean{Representation.charProjector,
        Representation.irreducibleCharacters,
        Representation.characterMultiplicity,
        Representation.deltaOperator}
    \leanok
    Let $\chi_U$ be the character of $U$, and let $\chi_1,\dots,\chi_r$ be the
    distinct characters of the irreducible subrepresentations of $U$, with
    $d_i=\chi_i(1)$ and multiplicities
    $m_i=\frac{1}{|G|}\sum_g\chi_U(g)\chi_i(g^{-1})$. The isotypic projectors
    are $P_i=\frac{d_i}{|G|}\sum_g\chi_i(g^{-1})U_g$, and
    \begin{align*}
        \Delta=\frac{1}{|G|}\sum_{i=1}^{r}\frac{d_i}{m_i}P_i
    \end{align*}
    \cite[Lemma~4.4]{Schuch2010PEPS}.
\end{definition}

\begin{lemma}[Isotypic projectors from characters]%
    \label{thm:peps_isotypic_projectors}
    \lean{Representation.sum_character_inv_smul_eq,
        Representation.exists_isInternal_isAtom,
        Representation.charProjector_apply_of_mem,
        Representation.irreducibleCharacters_finite,
        Representation.characterMultiplicity_eq_finrank,
        Representation.sum_charProjector_irreducibleCharacters,
        Representation.trace_comp_charProjector}
    \leanok
    \uses{def:peps_delta_operator}
    For irreducible representations $D^i$ and $D^j$ of $G$,
    $\sum_g\chi_i(g^{-1})D^j(g)$ equals $\frac{|G|}{d_i}\Id$ if $D^i$
    and $D^j$ are equivalent and $0$ otherwise. Consequently $P_i$ acts as
    the identity on every irreducible subrepresentation of $U$ with
    character $\chi_i$ and as zero on every other irreducible
    subrepresentation, the number $m_i$ is the dimension of the space of
    intertwining maps from $D^i$ to $U$, and
    \begin{align*}
        \sum_{i=1}^{r}P_i & =\Id, & \tr[U_kP_i] & =m_i\chi_i(k)
    \end{align*}
    \cite[Lemma~4.4]{Schuch2010PEPS}.
\end{lemma}

\begin{proof}\leanok
    The operator $\sum_g\chi_i(g^{-1})D^j(g)$ commutes with every $D^j(h)$,
    because characters are constant on conjugacy classes, so by Schur's lemma
    it is a multiple of the identity; its trace is
    $\sum_g\chi_i(g^{-1})\chi_j(g)$, which orthogonality of characters
    evaluates. By Maschke's theorem $U$ is a direct sum of irreducible
    subrepresentations. On each summand the projectors act as just
    described, which gives $\sum_iP_i=\Id$, and the trace of
    $U_kP_i$ is the sum of $\chi_i(k)$ over the summands with character
    $\chi_i$. Orthogonality of characters counts these summands as $m_i$.
\end{proof}

\begin{lemma}[Equality on irreducible summands]
    \label{thm:peps_irreducible_ext}
    \lean{Representation.linearMap_ext_on_irreducible}
    \leanok
    \uses{thm:peps_isotypic_projectors}
    Let $\rho$ be a finite-dimensional representation of a finite group
    over a field in which $|G|$ is nonzero. Two endomorphisms are equal
    if they agree on every irreducible subrepresentation of $\rho$.
    This is the decomposition principle behind the block calculations
    in \cite[Lemma~4.4]{Schuch2010PEPS}.
\end{lemma}

\begin{proof}\leanok
    Maschke's theorem expresses the representation space as an internal
    direct sum of irreducible subrepresentations. Agreement on these
    summands implies agreement on their span, which is the whole space.
\end{proof}

\begin{lemma}[Scalar calculus of character projectors]
    \label{thm:peps_weighted_character_projectors}
    \lean{Representation.sum_smul_charProjector_apply_of_mem,
        Representation.sum_smul_charProjector_commute,
        Representation.charProjector_mul_self}
    \leanok
    \uses{thm:peps_isotypic_projectors}
    Given complex weights $f(\chi)$ on the occurring irreducible characters,
    put $A_f=\sum_\chi f(\chi)P_\chi$. If $v$ lies in an irreducible summand
    with character $\psi$, then $A_fv=f(\psi)v$. In addition,
    $A_fU_g=U_gA_f$. Each projector $P_\chi$ is idempotent. These formulas
    express the isotypic decomposition used in
    \cite[Lemma~4.4]{Schuch2010PEPS}.
\end{lemma}

\begin{proof}\leanok
    On the summand with character $\psi$, the projector $P_\chi$ is the
    identity if $\chi=\psi$ and zero otherwise. Thus only the weight of
    $\psi$ survives. Each summand is preserved by the group action,
    which commutes with scalar multiplication. The irreducible summands
    span the whole representation space.
\end{proof}

\begin{lemma}[Block form of $\Delta$]%
    \label{thm:peps_delta_block_form}
    \lean{Representation.deltaOperator_apply_of_mem}
    \leanok
    \uses{def:peps_delta_operator}
    On every irreducible subrepresentation of $U$ with character $\chi_i$,
    $\Delta$ acts as the scalar $\frac{d_i}{m_i|G|}$. Thus, in a basis in which
    $U_g\cong\bigoplus_iD^i(g)\otimes\Id_{m_i}$,
    \begin{align*}
        \Delta\cong\frac{1}{|G|}\bigoplus_i\frac{d_i}{m_i}\,
        \Id_{d_i}\otimes\Id_{m_i}
    \end{align*}
    \cite[Lemma~4.4]{Schuch2010PEPS}.
\end{lemma}

\begin{proof}\leanok
    \uses{thm:peps_isotypic_projectors}
    Only the projector of $\chi_i$ acts nontrivially on such a
    subrepresentation, and it acts as the identity.
\end{proof}

\begin{lemma}[Expansion in the representation]%
    \label{thm:peps_delta_expansion}
    \lean{Representation.trace_comp_deltaOperator,
        Representation.sum_trace_comp_deltaOperator_smul,
        Representation.sum_trace_inv_comp_deltaOperator_comp_smul}
    \leanok
    \uses{def:peps_delta_operator}
    For every $g\in G$,
    \begin{align*}
        \sum_h\tr[U_{h^{-1}}\Delta U_g]\,U_h
        =\sum_h\tr[U_{gh^{-1}}\Delta]\,U_h=U_g .
    \end{align*}
    For a unitary representation $U_{h^{-1}}=U_h^\dagger$
    \cite[Lemma~4.4]{Schuch2010PEPS}.
\end{lemma}

\begin{proof}\leanok
    \uses{thm:peps_isotypic_projectors}
    By the trace formula for the projectors,
    $\tr[U_k\Delta]=\frac{1}{|G|}\sum_i d_i\chi_i(k)$, so
    $\sum_k\tr[U_{k^{-1}}\Delta]U_k=\sum_iP_i=\Id$. Substituting
    $h=kg$ gives
    $\sum_h\tr[U_{gh^{-1}}\Delta]U_h=\big(\sum_k\tr[U_{k^{-1}}\Delta]U_k\big)U_g
        =U_g$, and the first expression equals the second by cyclicity of
    the trace.
\end{proof}

\begin{definition}[Reciprocal weighted trace operator]
    \label{def:peps_delta_inverse}
    \lean{Representation.deltaOperatorInverse}
    \leanok
    \uses{def:peps_delta_operator}
    For the occurring irreducible characters, with dimensions $d_\chi$
    and multiplicities $m_\chi$, define
    \begin{align}
        \Delta^{-1}=\sum_\chi\frac{|G|m_\chi}{d_\chi}P_\chi.
        \label{eq:peps_delta_inverse}
    \end{align}
    These are the reciprocal coefficients of the operator $\Delta$
    of \cite[Lemma~4.4]{Schuch2010PEPS}.
\end{definition}

\begin{lemma}[Invertibility of $\Delta$]
    \label{thm:peps_delta_inverse}
    \lean{Representation.deltaOperatorInverse_apply_of_mem,
        Representation.deltaOperatorInverse_mul,
        Representation.deltaOperator_mul_inverse,
        Representation.isUnit_deltaOperator,
        Representation.deltaOperatorInverse_commute,
        Representation.deltaOperator_commute}
    \leanok
    \uses{def:peps_delta_inverse, thm:peps_weighted_character_projectors,
        thm:peps_delta_expansion}
    The operator in (\ref{eq:peps_delta_inverse}) is the two-sided inverse
    of $\Delta$. On a summand with character $\chi$, it acts by
    $|G|m_\chi/d_\chi$. Both $\Delta$ and its inverse commute with every
    group operator. These assertions hold without semi-regularity or
    unitarity, for every finite-dimensional complex representation of a
    finite group.
\end{lemma}

\begin{proof}\leanok
    The dimensions and multiplicities of the occurring summands are
    positive. On each irreducible summand the two coefficients multiply
    to one:
    $[d_\chi/(|G|m_\chi)][|G|m_\chi/d_\chi]=1$.
    Hence $\Delta^{-1}\Delta=\Id$ on their direct sum. In finite dimension
    a left inverse is also a right inverse. Commutation follows from
    the scalar calculus of the character projectors.
\end{proof}

\begin{lemma}[Adjoints, inversion, and unitary characters]
    \label{thm:peps_unitary_character_adjoint}
    \lean{LinearMap.trace_adjoint,
        Representation.character_inv_of_unitary}
    \leanok
    Let $E$ be a finite-dimensional complex inner-product space. For every
    linear endomorphism $A$ of $E$,
    \[
        \tr(A^\dagger)=\overline{\tr(A)}.
    \]
    Consequently, if $\rho$ is a unitary representation of a group $G$ on $E$,
    then its character satisfies
    \[
        \chi_\rho(g^{-1})=\overline{\chi_\rho(g)}\qquad(g\in G).
    \]
\end{lemma}

\begin{proof}\leanok
    Choose an orthonormal basis. The matrix of an adjoint is the conjugate
    transpose, whose trace is the conjugate trace. In a unitary representation,
    $\rho(g^{-1})=\rho(g)^\dagger$.
\end{proof}

\begin{lemma}[Positivity of $\Delta$]
    \label{thm:peps_delta_positive}
    \lean{Representation.adjoint_subrepresentation_of_unitary,
        Representation.charProjector_adjoint_of_unitary,
        Representation.isPositive_charProjector_of_unitary,
        Representation.isPositive_sum_smul_charProjector_of_nonneg,
        Representation.isPositive_deltaOperator_of_unitary,
        Representation.posDef_toMatrix_deltaOperator_of_unitary}
    \leanok
    \uses{thm:peps_delta_inverse, thm:peps_isotypic_projectors,
        thm:peps_unitary_character_adjoint}
    Let $U$ be unitary. Every invariant subrepresentation is unitary. For
    unitary representations, the character projector of any irreducible
    summand is self-adjoint and positive. More generally, every nonnegative
    weighted sum of the occurring character projectors is positive. In
    particular $\Delta$ is positive, and its matrix in every orthonormal basis
    is positive definite \cite[Lemma~4.4]{Schuch2010PEPS}.
\end{lemma}

\begin{proof}\leanok
    Unitarity gives $\chi(g^{-1})=\overline{\chi(g)}$. Hence the character
    projectors of the occurring summands are self-adjoint. They are also
    idempotent, so they are positive. The coefficients
    $d_\chi/(|G|m_\chi)$ are positive. Their weighted sum is therefore
    positive; its invertibility makes it positive definite.
\end{proof}

\begin{definition}[Fourth-root isotypic weights]%
    \label{def:peps_representation_theta}
    \lean{Representation.thetaOperator}
    \leanok
    \uses{def:peps_delta_operator}
    Define $\Theta=\sum_i(d_i/m_i)^{1/4}P_i$, with the positive real fourth
    root. For a representation in which every irreducible representation
    occurs once, as in \cite[Section~7]{Schuch2010PEPS}, this is the
    operator $\bigoplus_i d_i^{1/4}\Id_{d_i}$.
\end{definition}

\begin{lemma}[Action of the fourth-root weights]%
    \label{thm:peps_representation_theta_scalar}
    \lean{Representation.thetaOperator_apply_of_mem,
        Representation.thetaOperator_apply_of_multiplicity_one}
    \leanok
    \uses{def:peps_representation_theta, thm:peps_weighted_character_projectors}
    On an irreducible summand of character $\chi_i$, the operator $\Theta$
    is multiplication by $(d_i/m_i)^{1/4}$. If $m_i=1$, the scalar is
    $d_i^{1/4}$.
\end{lemma}

\begin{proof}\leanok
    Each isotypic projector is the identity on its own irreducible summands
    and zero on the others. The weighted sum retains precisely the coefficient
    of the given character.
\end{proof}

\begin{lemma}[Fourth power and symmetry of the weight operator]%
    \label{thm:peps_representation_theta_power}
    \lean{Representation.thetaOperator_pow_four,
        Representation.thetaOperator_commute}
    \leanok
    \uses{thm:peps_representation_theta_scalar, thm:peps_delta_block_form,
        thm:peps_weighted_character_projectors}
    The weight operator commutes with the representation and satisfies
    $\Theta^4=|G|\Delta$. These algebraic identities require neither unitarity
    nor semi-regularity. The factor $|G|$ corrects the printed equality
    in \cite[Section~7]{Schuch2010PEPS}, because its Lemma~4.4 includes
    $|G|^{-1}$ in the definition of $\Delta$; see
    \cite{gap:rmp_peps_quantum_double_g_isometry}.
\end{lemma}

\begin{proof}\leanok
    On every irreducible summand, the fourth power of the positive real root
    is $d_i/m_i$, while $\Delta$ acts by $d_i/(m_i|G|)$.
    Complete reducibility gives the operator identity. Each weighted isotypic
    sum commutes with the representation.
\end{proof}

\begin{lemma}[The square of the fourth-root weights]%
    \label{thm:peps_representation_theta_squared}
    \lean{Representation.thetaOperator_sq_apply_of_mem,
        Representation.thetaOperator_sq_apply_of_multiplicity_one}
    \leanok
    \uses{thm:peps_representation_theta_scalar}
    On an irreducible summand of dimension $d_i$ and multiplicity $m_i$,
    the operator $\Theta^2$ acts by $\sqrt{d_i/m_i}$. In particular, for
    multiplicity one it acts by $\sqrt{d_i}$, the weight appearing in
    \cite[Section~7]{Schuch2010PEPS}.
\end{lemma}

\begin{proof}\leanok
    The operator $\Theta$ acts by $(d_i/m_i)^{1/4}$ on the summand.
    Applying it twice squares this positive real fourth root and gives
    $\sqrt{d_i/m_i}$. Set $m_i=1$ for the second assertion.
\end{proof}

\begin{lemma}[Translated irreducible character pairings]
    \label{thm:peps_translated_character_pairing}
    \lean{Representation.sum_character_inv_mul_character_translate,
        Representation.sum_character_inv_mul_character_translate_eq_zero}
    \leanok
    \uses{thm:peps_isotypic_projectors}
    Let $\sigma,\tau$ be finite-dimensional irreducible complex
    representations of $G$, and write $d_\sigma=\dim\sigma$. For all
    $a,b\in G$,
    \[
        \sum_{g\in G}\chi_\sigma((ag)^{-1})\chi_\tau(bg)
        =\begin{cases}
            \dfrac{|G|}{d_\sigma}\chi_\tau(ba^{-1}), & \sigma\cong\tau,     \\
            0,                                       & \sigma\not\cong\tau.
        \end{cases}
    \]
    In particular, if $\chi_\sigma\ne\chi_\tau$, every such translated pairing
    vanishes. These identities are used in \cite{Schuch2010PEPS}, lines
    2464--2486 of the source.
\end{lemma}

\begin{proof}\leanok
    Change variables from $g$ to $ag$. Apply character-weighted Schur
    orthogonality to $\sum_g\chi_\sigma(g^{-1})\tau(g)$, multiply by
    $\tau(ba^{-1})$, and take the trace. The sum is $|G|/d_\sigma$ times the
    identity when the irreducible representations are equivalent, and zero
    otherwise. Equality of irreducible characters is equivalent to equivalence
    of their representations.
\end{proof}

\begin{lemma}[Orthogonality and norm of translated characters]
    \label{thm:peps_translated_character_inner_product}
    \lean{Representation.sum_star_character_mul_character_translate,
        Representation.sum_star_character_mul_character_translate_eq_zero,
        Representation.sum_star_character_mul_character_self}
    \leanok
    \uses{thm:peps_translated_character_pairing,
        thm:peps_unitary_character_adjoint}
    Under the preceding hypotheses, suppose $\sigma$ is unitary. Then
    \[
        \sum_{g\in G}\overline{\chi_\sigma(ag)}\chi_\tau(bg)
        =\begin{cases}
            \dfrac{|G|}{d_\sigma}\chi_\tau(ba^{-1}), & \sigma\cong\tau,     \\
            0,                                       & \sigma\not\cong\tau.
        \end{cases}
    \]
    Thus translates of distinct irreducible characters are orthogonal,
    independently of $a$ and $b$, and each translate of the character of a
    unitary irreducible representation has squared norm
    \[
        \sum_{g\in G}|\chi_\sigma(ag)|^2=|G|.
    \]
    Translates of the same irreducible character need not be orthogonal.
\end{lemma}

\begin{proof}\leanok
    Replace the conjugated character by its value at the inverse group
    element. For the norm, put $\tau=\sigma$ and $b=a$; the factor
    $\chi_\sigma(1)=d_\sigma$ cancels the irreducible dimension.
\end{proof}
